\begin{tikzpicture}[font=\small,>=stealth]
\draw[rounded corners=3pt,gray!45] (0,0) rectangle (14.2,12.6);
\node at (7.1,12.05) {再放回椭圆：从“横向位移”看清“到准线距离”};\begin{scope}[shift={(4.2,6.25)},scale=0.62]
\draw[->,gray] (-5.6,0)--(5.7,0);
\draw[teal!75!black,thick] (0,0) ellipse (5 and 3);
\draw[blue!75!black,very thick] (4,0)--(4.6428571429,1.1134612334);\draw[orange!85!black,very thick] (4,0)--(2.5000000000,-2.5980762114);
\fill (4.6428571429,1.1134612334) circle(2pt) node[above=8pt] {$A$};\fill (2.5000000000,-2.5980762114) circle(2pt) node[below=8pt] {$B$};
\fill (4,0) circle(2pt);\node at (2.7,-.9) {$F$};\draw[gray,thin] (3.3,-.65)--(3.95,-.03);

\draw[violet,dashed] (6.25,-4.2)--(6.25,4.6);
\node[violet,above=7pt] at (6.25,4.6) {$\ell$};
\draw[blue!65!black,dashed] (4.6428571429,1.1134612334)--(6.25,1.1134612334371354);
\draw[orange!85!black,dashed] (2.5000000000,-2.5980762114)--(6.25,-2.598076211353316);
\fill (6.25,1.1134612334371354) circle(1.4pt) node[right=6pt] {$M$};
\fill (6.25,-2.598076211353316) circle(1.4pt) node[right=6pt] {$N$};
\end{scope}
\node[anchor=west,align=left] at (9.5,7.6)
 {$\ell:x=a^2/c>a$\\[12pt]
 $x_A=c+r\cos\theta$\\[12pt]$x_B=c-s\cos\theta$\\[15pt]
 $AF=e\,AM$\\[12pt]$BF=e\,BN$\\[15pt]
 正向、反向的投影\\[8pt]符号正好相反};
\node at (7.1,2.8) {$\textcolor{blue!75!black}{AM=\dfrac{b^2}{c}-r\cos\theta}$};
\node at (7.1,1.65) {$\textcolor{orange!85!black}{BN=\dfrac{b^2}{c}+s\cos\theta}$};
\node at (7.1,.55) {先保留投影结构，再用第二定义；准线在整个椭圆的右侧。};

\end{tikzpicture}
